\documentclass[12pt]{article}

% chktex-file 1
% chktex-file 3

\usepackage[utf8]{inputenc}
\usepackage[T1]{fontenc}
\usepackage{amsmath, amssymb}
\usepackage{graphicx}
\usepackage{tikz}
\usepackage{geometry}
\usepackage{booktabs}

\geometry{margin=2.5cm}
\begin{document}

\section{Parte 1: Por transformada inversa}

La función de densidad proporcionada es:

\[
f(x)=
\begin{cases}
\dfrac{(x-3)^2}{18}{}, & 0 \leq x \leq 6 \\
0, & x < 0 \text{ o } x > 6
\end{cases}
\]

El objetivo consiste en generar valores aleatorios mediante el método de transformada inversa y verificar que los valores obtenidos pertenezcan al intervalo $[0,6]$.

\subsection{Representación gráfica referencial}

\begin{center}
\begin{tikzpicture}[scale=1.2]

% Ejes
\draw[->] (0,0) -- (8,0) node[right] {$x$};
\draw[->] (0,0) -- (0,4) node[above] {$f(x)$};

% Triángulo
\draw (1.5,0) -- (4,3) -- (6.5,0);

% Línea punteada
\draw[dashed] (4,0) -- (4,3);

% Etiquetas
\node[below] at (1.5,0) {$a$};
\node[below] at (4,0) {$b$};
\node[below] at (6.5,0) {$c$};

\end{tikzpicture}
\end{center}

\section{Parte 2: Aplicaciones de Simulación}

\subsection{Problema 1}

Una cierta compañía posee un gran número de máquinas en uso.

El tiempo que dura en operación cada máquina sigue la siguiente distribución:

\begin{center}
\begin{tabular}{cc}
\toprule
Tiempo entre descomposturas (horas) & Probabilidad \\
\midrule
6--8   & 0.10 \\
8--10  & 0.15 \\
10--12 & 0.24 \\
12--14 & 0.26 \\
16--18 & 0.18 \\
18--20 & 0.07 \\
\bottomrule
\end{tabular}
\end{center}

El tiempo de reparación sigue:

\begin{center}
\begin{tabular}{cc}
\toprule
Tiempo de reparación (horas) & Probabilidad \\
\midrule
2--4   & 0.15 \\
4--6   & 0.25 \\
6--8   & 0.30 \\
8--10  & 0.20 \\
10--12 & 0.10 \\
\bottomrule
\end{tabular}
\end{center}

Si el costo de tener una máquina ociosa es de \$500 por hora y el salario del operario es de \$50 por hora, se debe determinar cuántas máquinas asignar a cada mecánico.

\subsection{Problema 2}

Una cadena de supermercados es abastecida por un almacén central.

Los camiones llegan según un proceso Poisson con razón media de 2 camiones por hora.

Los tiempos de descarga dependen del número de trabajadores:

\begin{center}
\begin{tabular}{cc}
\toprule
Número de trabajadores & Tiempo de descarga (min) \\
\midrule
3 & Uniforme (20,30) \\
4 & Uniforme (15,25) \\
5 & Uniforme (10,20) \\
6 & Uniforme (5,15) \\
\bottomrule
\end{tabular}
\end{center}

Cada trabajador recibe \$25 por hora durante un turno nocturno de 8 horas.

El costo de espera de un camión es de \$50 por hora.

El objetivo es determinar el tamaño óptimo del equipo.

\end{document}